\documentclass[twocolumn]{article} % 双栏排版，符合多数物理通报期刊风格
% 中文支持
\usepackage[UTF8]{ctex}            % 支持中文
\usepackage{amsmath, amssymb, amsthm}      % 数学符号与公式
\usepackage{graphicx}              % 插入图片
\usepackage{caption}               % 图表标题
\usepackage{float}                 % 图表位置控制
\usepackage{geometry}              % 页面设置
\usepackage{fancyhdr}              % 页眉页脚
\usepackage{abstract}              % 摘要环境优化
\usepackage{indentfirst}           % 首段缩进
\usepackage{setspace}              % 行距设置
\usepackage{hyperref}              % 超链接
\usepackage{dutchcal}
\usepackage{bm}
\usepackage{fmtcount}
\usepackage{mathtools}
\usepackage{booktabs}
\usepackage{array}
\usepackage{tikz}
\usepackage[european resistors]{circuitikz}
\usepackage{siunitx}
\usepackage{subcaption}
\usepackage{pgfplots}
\usepackage{mathrsfs}
\pgfplotsset{compat=1.18}
\newcommand{\di}{\mathrm{d}}
\newcommand{\ii}{\mathrm{i}}
\newcommand{\e}[1]{\mathrm{e}^{#1}}
\newcommand{\f}[1]{\mathscr{F}\left[#1\right]}
\newcommand{\fo}[1]{\mathscr{F}^{-1}\left[#1\right]}
\ctikzset{
%    label/align = rotate,
    bipoles/length = 1cm,      % 元件长度
    resistors/scale = 0.8,     % 电阻符号缩放
    capacitors/scale = 0.8,    % 电容符号缩放
}
% 页面设置（符合中文期刊习惯）
\geometry{
  a4paper,
  left=1.0cm,
  right=1.0cm,
  top=2.5cm,
  bottom=2.5cm,
  columnsep=0.8cm                 % 双栏间距
}
% 行距
\setstretch{1.2}
% 标题格式
\title{\huge \bfseries 网格电阻的Fourier分析}
\author{
  朱克拉$^{1}$，\quad 王思晗$^{1}$ \\
  \small (1.中国科学技术大学工程科学学院，合肥)
}
\date{}
% 重定义摘要环境为单栏，跨双栏
\renewenvironment{abstract}{\small
  \textbf{摘\ \ 要：}
}{}
\newenvironment{keyword}{\small
  \textbf{关键词：}  
}{}
% 首行缩进
\setlength{\parindent}{2em}
% 图表标题格式
\captionsetup[figure]{labelformat=default, labelsep=quad, font=small}
\captionsetup[table]{labelformat=default, labelsep=quad, font=small}
% 公式编号格式
%\renewcommand{\theequation}{\arabic{equation}}
%%%%%%%%%%%%%%%%%%%%%%%%%%%%%%%%%%%%%%%%%%%%%%%%%%%%%%%%%%%%%%%%%5
\begin{document}

% 单栏输出标题、作者、摘要、关键词（跨双栏）
\twocolumn[
  \begin{@twocolumnfalse}
    \maketitle
    \begin{abstract}
      本文提供一份适用于《物理学通报》类期刊的LaTeX论文模板。展示了中文论文的基本结构，
      包括标题、作者信息、摘要、关键词、公式推导、图表引用以及参考文献的格式。
      用户可在此基础上直接替换内容，快速完成稿件排版。
    \end{abstract}


    \begin{keyword}
      物理学通报；LaTeX模板；双栏排版；物理公式
    \end{keyword}
    \vspace{0.5cm}
  \end{@twocolumnfalse}
]

% 正文开始（双栏）
\section{引言}

《物理学通报》是国内物理教学与科研交流的重要期刊之一，涵盖物理理论研究、实验方法、
教学探讨等内容。使用LaTeX可以生成高质量、符合出版规范的文档，尤其适合包含复杂公式
和多图表的论文。本文提供一个基础模板，包含常见的排版要素。

\section{理论基础}
设电流分布
\begin{align}
  I(x,y)=\begin{dcases}
    I,&x=y=0,\\
    -I,&x=m,y=n,\\
    0,&\text{other}.
  \end{dcases}
\end{align}

$V(x,y)$是每个点的电势，有KCL
\begin{align*}
  I(x,y)&=\dfrac{1}{r}\left[V(x,y)-V(x-1,y)\right]\\
        &+\dfrac{1}{r}\left[V(x,y)-V(x+1,y)\right]\\
        &+\dfrac{1}{r}\left[V(x,y)-V(x,y-1)\right]\\
        &+\dfrac{1}{r}\left[V(x,y)-V(x,y+1)\right]
\end{align*}

总结起来
\begin{equation}
  V(x,y)-\mathscr{L}V(x,y)=\dfrac{r}{4}I(x,y)
\end{equation}

其中，
\begin{align*}
  \mathscr{L}V(x,y)&=\dfrac{1}{4}\left[V(x-1,y)+V(x+1,y)\right.\\
        &\left.+V(x,y-1)+V(x,y+1)\right]
\end{align*}

所要求的电阻值
\begin{equation}
  R(m,n)=\dfrac{1}{I}\left[V(0,0)-V(m,n)\right]
\end{equation}

设$V(x,y)=\f{F(\omega_1,\omega_2)}$，于是有
\begin{equation}
  F(\omega_1,\omega_2)=\fo{V(x,y)}
\end{equation}

考虑$I(x,y)$的Fourier逆变换，
\begin{align*}
  \fo{I(x,y)}&=\sum_{x=-\infty}^{\infty}\sum_{y=-\infty}^{\infty}I(x,y)\e{\ii(x\omega_1+y\omega_2)}\\
  &=I\left[1-\e{\ii(m\omega_1+n\omega_2)}\right]
\end{align*}

利用Fourier变换的频移性，不难验证
\begin{equation}
  \fo{\mathscr{L}V(x,y)}=F(\omega_1,\omega_2)\dfrac{\cos\omega_1+\cos\omega_2}{2}
\end{equation}

对于式(?)两端进行Fourier逆变换，有
\begin{equation}
  F(\omega_1,\omega_2)=\dfrac{Ir}{2}\dfrac{1-\e{\ii(m\omega_1+n\omega_2)}}{2-(\cos\omega_1+\cos\omega_2)}
\end{equation}

于是，$V(x,y)$只需对$F(\omega_1,\omega_2)$进行Fourier变换
\begin{align*}
  &V(x,y)=\f{F(\omega_1,\omega_2)}\\
  &=\dfrac{Ir}{8\pi^2}\iint_{\Omega}\dfrac{\left[1-\e{\ii(m\omega_1+n\omega_2)}\right]\e{-\ii(\omega_1x+\omega_2y)}}{2-(\cos\omega_1+\cos\omega_2)}\di\omega_1\di\omega_2
\end{align*}
其中，$\Omega=[-\pi,\pi]^2$。

考虑积分区间的对称性，虚部消失，于是有
\begin{align}
  \begin{dcases}
    V(0,0)=\dfrac{Ir}{8\pi^2}\iint_{\Omega}\dfrac{1-\cos{(m\omega_1+n\omega_2)}}{2-(\cos\omega_1+\cos\omega_2)}\di\omega_1\di\omega_2\\
    V(m,n)=-\dfrac{Ir}{8\pi^2}\iint_{\Omega}\dfrac{1-\cos{(m\omega_1+n\omega_2)}}{2-(\cos\omega_1+\cos\omega_2)}\di\omega_1\di\omega_2
  \end{dcases}
\end{align}

得出电阻的积分解为
\begin{equation}
  R_{mn}=\dfrac{r}{4\pi^2}\iint_{\Omega}\dfrac{1-\cos{(m\omega_1+n\omega_2)}}{2-(\cos\omega_1+\cos\omega_2)}\di\omega_1\di\omega_2
\end{equation}
但是这个积分式不适合计算，下面计算其级数形式。

记积分式为
\begin{equation}
  I(m,n)=\iint_{\Omega}\dfrac{1-\cos{(m\omega_1+n\omega_2)}}{2-(\cos\omega_1+\cos\omega_2)}\di\omega_1\di\omega_2
\end{equation}

作变换$u=\dfrac{\omega_1+\omega_2}{2},v=\dfrac{\omega_1-\omega_2}{2}$，又有
\begin{equation}
  \di\omega_1\di\omega_2=\left|\dfrac{\partial(\omega_1,\omega_2)}{\partial(u,v)}\right|\di u\di v=2\di u\di v
\end{equation}
\begin{equation}
  \Omega\mapsto\Omega'
\end{equation}
积分区域变换为$\Omega'=\left\{(u,v)||u|+|v|\leqslant\pi\right\}$，令$a=m+n,b=m-n$。


则积分式写作
\begin{align*}
  I(a,b)&=\iint_{\Omega'}\dfrac{1-\cos{(au+bv)}}{1-\cos u\cos v}\di u\di v\\
  &=\iint_{\Omega'}(1-\cos{(au+bv)})\sum_{k=0}^{\infty}\cos^ku\cos^kv\di u\di v\\
  &=I_1-I_2
\end{align*}

拆成的两个积分式
\begin{equation}
  I_1=\sum_{k=0}^{\infty}\iint_{\Omega'}\cos^ku\cos^kv\di u\di v
\end{equation}
\begin{equation}
  I_2=\sum_{k=0}^{\infty}\iint_{\Omega'}\cos(au+bv)\cos^ku\cos^kv\di u\di v
\end{equation}

其中，
\begin{align*}
  I_1&=2\sum_{k=0}^{\infty}\iint_{[0,\pi]^2}\cos^ku\cos^kv\di u\di v\\
  &=2\sum_{k=0}^{\infty}\int_0^{\pi}\cos^ku\di u\int_0^{\pi}\cos^kv\di v\\
  &=2\sum_{k=\text{偶数}}^{\infty}\left(2\times\dfrac{(k-1)!!}{k!!}\dfrac{\pi}{2}\right)^2\\
  &=2\pi^2\sum_{k=0}^{\infty}\left(\dfrac{(2k-1)!!}{(2k)!!}\right)^2\\
  &=2\pi^2\sum_{k=0}^{\infty}\dfrac{1}{2^{4k}}\dfrac{(2k)!^2}{k!^4}
\end{align*}

再看$I_2$，
\begin{align*}
  I_2&=\sum_{k=0}^{\infty}\iint_{\Omega'}\cos(au+bv)\cos^ku\cos^kv\di u\di v\\
  &=\sum_{k=0}^{\infty}\iint_{\Omega'}\cos au\cos bv\cos^ku\cos^kv\di u\di v\\
  &=2\sum_{k=0}^{\infty}\left(\int_0^\pi\cos au\cos^ku\di u\right)\left(\int_0^\pi\cos bv\cos^kv\di v\right)
\end{align*}

记
\begin{equation}
  J(a)=\dfrac{1}{2}\int_{-\pi}^\pi\cos au\cos^ku\di u
\end{equation}

令$z=\e{\ii u}$，则有
\begin{equation}
  \di z=\ii\e{\ii u}\di u=\ii z\di u
\end{equation}
\begin{align*}
  J(a)&=\dfrac{1}{2\ii}\oint_{|z|\leqslant1}\dfrac{z^a+z^{-a}}{2}\dfrac{\left(z+z^{-1}\right)^k}{2^k}z^{-1}\di z\\
  &=\dfrac{1}{2^{k+2}\ii}\oint_{|z|\leqslant1}\dfrac{\left(z^{2a}+1\right)}{z^{a+k+1}}\left(z^2+1\right)^k\di z\\
  &=\dfrac{1}{2^{k+2}\ii}\oint_{|z|\leqslant1}\sum_{t=0}^{k}\binom{k}{t}z^{2t}\dfrac{\left(z^{2a}+1\right)}{z^{a+k+1}}\di z\\
  &=\dfrac{-\ii}{2^{k+2}}\sum_{t=0}^{k}\binom{k}{t}\oint_{|z|\leqslant1}\left(\dfrac{1}{z^{k+a+1-2t}}+\dfrac{1}{z^{k-a+1-2t}}\right)\di z\\
\end{align*}

利用留数定理，只需算被积函数洛朗展开$z^{-1}$系数
\begin{align*}
  J(a)&=\dfrac{\pi}{2^{k+1}}\left[\binom{k}{\frac{k+a}{2}}+\binom{k}{\frac{k-a}{2}}\right]\\
  &=\dfrac{\pi}{2^k}\dfrac{k!}{\left(\frac{k+a}{2}\right)!\left(\frac{k-a}{2}\right)!}
\end{align*}

考虑$J(a)$关于$k$的取值，
\begin{equation}
  J(a)=
  \begin{dcases}
    \dfrac{\pi}{2^k}\dfrac{k!}{\left(\frac{k+a}{2}\right)!\left(\frac{k-a}{2}\right)!},&|a|\leqslant k,k\equiv a(\text{mod}2)\\
    0,&\text{other}.
  \end{dcases}
\end{equation}

同理于$J(a)$，
\begin{equation}
  J(b)=
  \begin{dcases}
    \dfrac{\pi}{2^k}\dfrac{k!}{\left(\frac{k+b}{2}\right)!\left(\frac{k-b}{2}\right)!},&|b|\leqslant k,k\equiv b(\text{mod}2)\\
    0,&\text{other}.
  \end{dcases}
\end{equation}

考虑$J(a)J(b)$：
\begin{enumerate}
  \item $a,b$必然同奇偶；
  \item 三角不等式$|a|\geqslant|b|$。
\end{enumerate}
故$I_2$的求和中只有$k=a+2p,p=0,1,2,\cdots$对求和有贡献。所以，$I_2$中关于$k$的求和换成对$p$的求和
\begin{align}
  I_2=\sum_{p=0}^{\infty}\dfrac{2\pi^2}{2^{2a+4p}}\dfrac{(a+2p)!}{p!(a+p)!}\dfrac{(a+2p)!}{(\frac{a-b}{2}+p)!(\frac{a+b}{2}+p)!}
\end{align}

注意$a,b$的定义式，并将$p$替换回$k$，又有
\begin{equation}
  I_2=\dfrac{2\pi^2}{2^{2(m+n)}}\sum_{k=0}^{\infty}\dfrac{1}{2^{4k}}\dfrac{(m+n+2k)!}{k!(m+k)!(n+k)!(m+n+k)!}
\end{equation}

结合$I_1,I_2$的表达式，得出积分$I$的级数为
\begin{align*}
  I&=I_1-I_2\\
  &=2\pi^2\sum_{k=0}^{\infty}\dfrac{1}{2^{4k}}\left[\left(\dfrac{(2k)!}{k!k!}\right)^2\right.\\
  &\left.-\dfrac{(m+n+2k)!^2}{2^{2(m+n)}k!(m+k)!(n+k)!(m+n+k)!}\right]
\end{align*}

最终得出的电阻公式为
\begin{align*}
  &R_{mn}=\dfrac{r}{2}\sum_{k=0}^{\infty}\dfrac{1}{2^{4k}}\left[\left(\dfrac{(2k)!}{k!k!}\right)^2\right.\\
  &\left.-\dfrac{(m+n+2k)!^2}{2^{2(m+n)}k!(m+k)!(n+k)!(m+n+k)!}\right]
\end{align*}




\section{实验与讨论}

图\ref{fig:example}是示例图片。如图可见，变化趋势符合理论预期。

\begin{figure}[H]
  \centering
  \includegraphics[width=0.8\columnwidth]{example-image} % 请替换为实际图片
  \caption{示例图片（请替换为实际图像文件）}
  \label{fig:example}
\end{figure}

表\ref{tab:data}给出了部分测量数据。

\begin{table}[htbp]
  \centering
  \caption{测量数据示例}
  \label{tab:data}
  \begin{tabular}{|c|c|c|}
    \hline
    序号 & 温度 (K) & 电阻 ($\Omega$) \\
    \hline
    1 & 300 & 22.4 \\
    2 & 350 & 26.7 \\
    3 & 400 & 31.2 \\
    \hline
  \end{tabular}
\end{table}

\section{结论}

本文给出的LaTeX模板可直接用于撰写《物理学通报》风格的论文。用户需注意：

\begin{itemize}
  \item 所有图片文件放在同一目录，并用正确的文件名替换示例中的占位符；
  \item 参考文献请使用BibTeX或手动按顺序编号；
  \item 双栏排版下，长表格或宽图可用\texttt{table*}或\texttt{figure*}环境使之跨栏；
  \item 如果需要补充中文摘要和英文摘要，可以使用\texttt{abstract}和\texttt{keywords}环境。
\end{itemize}

\section*{参考文献}

[1]舒幼生,胡望雨,陈秉乾.物理学难题集萃. 下册[M].中国科学技术大学出版社,2014:611-616.

[2] Einstein A. Zur Elektrodynamik bewegter Körper[J]. Annalen der Physik, 1905, 17: 891-921.

[3] 张三， 李四. 量子纠缠的物理教学探讨[J]. 物理通报， 2023, (5): 45-49.

\end{document}